% ============================================================================
% DHG508 Final Design Paper & Post-Mortem — zhanguo_qunxiong (Godot build)
% OVERLEAF SETUP:
%   1. Upload this paper/ folder (or dhg508-overleaf-template.zip)
%   2. Menu (top left) -> Compiler -> XeLaTeX
%   3. Click Recompile
% ============================================================================

\documentclass[12pt,a4paper]{article}

\usepackage{fontspec}
\newcommand{\MainLatinFont}{Times New Roman}
\newcommand{\FallbackLatinFont}{TeX Gyre Termes}
\newcommand{\MainCJKFont}{Noto Serif CJK SC}
\newcommand{\FallbackCJKFont}{FandolSong-Regular.otf}

\IfFontExistsTF{\MainLatinFont}
  {\setmainfont{\MainLatinFont}}
  {\setmainfont{\FallbackLatinFont}}

\usepackage{xeCJK}
\IfFontExistsTF{\MainCJKFont}
  {\setCJKmainfont{\MainCJKFont}}
  {\IfFontExistsTF{Songti SC}
    {\setCJKmainfont{Songti SC}}
    {\IfFontExistsTF{SimSun}
      {\setCJKmainfont{SimSun}}
      {\setCJKmainfont{\FallbackCJKFont}}}}

\usepackage[a4paper, margin=1in]{geometry}
\usepackage{setspace}
\usepackage{booktabs}
\usepackage{array}
\usepackage{graphicx}
\usepackage{csquotes}

% Chicago Manual of Style, Notes-Bibliography (BibLaTeX + Biber)
\usepackage[notes, backend=biber, autolang=other, isbn=false]{biblatex-chicago}
\addbibresource{references.bib}

\renewcommand{\footnotesize}{\fontsize{10}{12}\selectfont}
\usepackage[hidelinks]{hyperref}

\title{\textbf{知識即戰力：史料上臉的戰國群雄策略游戲}\\[0.3em]
\large Knowledge as Combat Power: A Warring States Strategy Game with Its Historiography on Its Face\\[0.2em]
\large Design Rationale, Evidence Boundaries, and Post-Mortem}
\author{James\\Lingnan University $\cdot$ DHG508 Critical Digital History\\
\footnotesize Live: \url{https://zhanguo-godot.zhanguo-godot.workers.dev}\\
\footnotesize Repo: \url{https://github.com/wangsizhangyang-cloud/dhg508}}
\date{25 November 2026}

\begin{document}
\maketitle

\begin{abstract}
\noindent
This paper presents \emph{Zhanguo Qunxiong} (戰國群雄), a turn-based grand
strategy simulation of the Warring States, shipped as a Godot~4 build on
Cloudflare Workers. The player picks one of the seven great states and
competes to unify China through a HOI4-style loop of focus trees, domestic
policy, diplomacy, and formation-based battles---with every number traceable
to a source. The game's data layer is extracted from the \emph{Shiji}, the
\emph{Guoyu}, the \emph{Xunzi}, and the legal-military corpus; every claim is
tagged \emph{据} (documented), \emph{推断} (inferred), or \emph{设}
(invented); and the focus-tree and law nodes quote their sources inside the
game itself. Seventeen headless Godot self-check scripts and Playwright
end-to-end tests against two deployed URLs keep the build honest. The main
takeaway: a strategy game can carry its historiography on its face---sources
on the nodes, limits in the table---without ceasing to be a game.
\end{abstract}

\onehalfspacing

\section{Historical Question and Medium}

The project asks two questions at once. The historical question: what made
the Warring States---a world of fanatical reform, mass infantry, and
diplomatic improvisation---cohere as a system, and can a player feel that
system work? The design question follows: can a grand strategy game be built
so that \emph{every} number on screen is either attested, honestly inferred,
or openly invented---without the game collapsing under its own footnotes?

An interactive simulation is the right medium because the answer to the
historical question is systemic. Reading that ``Qi's hired retainers are no
match for Wei's \emph{wuzu}, and Wei's \emph{wuzu} are no match for Qin's
\emph{ruishi}''\autocite{xunzi} states a hierarchy; only playing Qin against
Wei makes the player \emph{pay} for the difference---in upkeep, in recruitment
cost, in which provinces are worth taking first. The game operationalises the
course's central claim that historical knowledge should be a spendable asset:
reading the sources in the codex is not homework beside the game, it is how
you learn which reforms to rush and which provinces to want.

The shipped build is the Godot~4 version: a full conquest game on an
87-province map, deployable as a WebGL build that anyone can open in a
browser. A web-rules prototype preceded it and remains as the reference
implementation; a tutorial episode and an LLM agent built in that prototype
were left out of the shipped build (Section~\ref{sec:postmortem}).

\section{Sources and Scholarly Conversation}

The primary corpus is the \emph{Shiji} of Sima Qian---master text for
persons, battle records, and state narratives---cross-checked against the
\emph{Guoyu}\autocite{guoyu} and, for military organisation, the
\emph{Xunzi}'s chapter ``Yi bing'' and the legal-military texts. The limits
are stated plainly: these are received texts, not excavated ones; the
1215 battle records in the database were verified for internal integrity and
spot-checked against paralleled texts, not re-edited line by line; and
Chinese primary-source citations inside the game use the received
chapter-and-passage convention, which some editions number differently.

Three scholarly works anchor the design conversation. First, Harnes's scoping
review of forty-one empirical studies concludes that historical empathy is a
cognitive--affective composite in which \emph{substantive knowledge} is
fundamental and contextualisation and evidence-work are the real
bottlenecks.\autocite{harnes2025empathy} This licensed the game's central
mechanism: knowledge is a settlement term---it unlocks reforms, gates failed
paths, and enters the combat formula---rather than a prerequisite we merely
hoped players would bring. It also warned against first-person role-play,
which heightens presentism; the shipped game accordingly has no avatar to be.

Second, Park et al.'s generative-agents architecture showed how LLM agents
can sustain memory, perception, and action---and, crucially, how retrieval
fails and fabricates.\autocite{park2023generative} The project built a King
You agent on the narrowest viable slice of that design: a hand-curated memory
whitelist drawn only from events attested in the ``Zhou benji,'' a bounded
perception window, and a single-action JSON output contract checked by the
world engine before anything is shown. The agent may improvise in register;
it may not improvise in fact.

Third, Zhang's re-dating of the Western Zhou chronology demonstrates that
771~BCE is not an undisputed anchor: the accession chronology can be
reconstructed differently, and 771 is itself contestable.\autocite{zhang2023chronology}
Rather than hide this, the game's annals and historical codex carry the
dispute as content: a player who digests the source-criticism dossier learns
that the calendar they are playing on is itself an argument.

\section{Architecture and Boundaries}

The system separates layers so that machines, not only humans, can verify the
historical content. A \textbf{data layer} (JSON, exported from a shared
SQLite source by a build script) holds 190 generals with 1215 attested battle
records, each carrying \texttt{original\_text} and a source citation; 15 unit
types and 13 counter relations converted from textual comparisons; 87
provinces anchored on historically attested cities digitised from Tan
Qixiang's historical atlas;\autocite{tan1982atlas} and an officials roster. A
\textbf{rules layer} in GDScript resolves the turn loop---revenue, training,
movement, siege, events---through combat formulas that consume unit
composition, counters, terrain, fortification, commanders, and four
formations (line, wedge, square, arrow). A \textbf{strategy layer} carries
the HOI4-style focus tree, whose branches follow the schools of the age
(\emph{fashu}, warfare, economy, infrastructure), each node quoting its
source---the root of Qin's tree cites the \emph{Shiji} on Xian's ambition,
``Agronomy and War'' cites the \emph{Shang jun shu}. A \textbf{diplomacy
layer} handles treaties, mediation, and negotiation, with minor states acting
on interest rather than scripted loyalty.

Deployment is a static Godot~4 WebGL export behind a small Cloudflare Worker:
the worker 302-redirects \texttt{/} to a no-store boot page (to bypass stale
edge cache), then streams \texttt{index3.wasm} and \texttt{index3.pck},
compressed at the edge according to \texttt{Accept-Encoding}. A GitHub Pages
mirror serves the uncompressed originals. Epistemic cutoffs are temporal as
well as structural: nothing in the data layer speaks of events after the
unification of 221~BCE, and the game's invented elements---formulas,
formations' numeric weights, map geometry---are tagged \emph{设} in the
evidence table below. The tag taxonomy follows the course convention:
\emph{据} documented, \emph{推断} inferred, \emph{设} invented.

One scope decision belongs here rather than in a footnote. The six-lesson
tutorial and the LLM King You agent were built and verified in the web
prototype, but the final scope freeze shipped the Godot build \emph{without}
them. The design lessons carry over---memory whitelisting, output validation,
and the presentism-driven refusal to give the player a historical avatar to
inhabit---but the public URL is, deliberately, a game in which the past is
met through events, annals, and cited reforms, not through a chat window.

\section{Evidence Table}

Table~\ref{tab:evidence} is generated from \texttt{evidence.csv} by
\texttt{make\_evidence\_table.py}; the LaTeX table is never hand-edited.

\begin{table}[htbp]
  \centering
  \scriptsize
  \singlespacing
  \renewcommand{\arraystretch}{2.2}
  \setlength{\extrarowheight}{2pt}
  % Auto-generated from evidence.csv by make_evidence_table.py — edit the CSV, not this file.
\begin{tabular}{p{0.27\textwidth}p{0.20\textwidth}p{0.13\textwidth}p{0.08\textwidth}p{0.22\textwidth}}
  \toprule
  \textbf{Claim / Element} & \textbf{Source / File} & \textbf{Locator} & \textbf{Level} & \textbf{Boundary / Limit} \\
  \midrule
  Deposition of Queen Shen and Crown Prince Yijiu; Marquis of Shen allies with Zeng and the Quanrong; You is killed at the foot of Mount Li & Sima Qian, Shiji, 'Zhou benji' and 'Qin benji' & Shiji 4.149; Shiji 5.179 & 据 & Received text accepted for teaching; variant accounts exist \\
  The beacon-fire episode: You lights the beacons but the lords no longer come & Shiji, 'Zhou benji' & Shiji 4.147-148 & 据 & Received-text narrative; later scholarship questions its historicity \\
  Duke Xiang of Qin escorts Prince Ping east and is enfeoffed as a hou west of the Qi river & Shiji, 'Qin benji' & Shiji 5.179 & 据 & Boundary of the grant ('Qi between the rivers') is summarized from the text \\
  Bo Yangfu and Shi Bo foresee the fall of Western Zhou & Guoyu, 'Zhou yu xia' and 'Zheng yu' & Guoyu 1; Guoyu 16 & 据 & Dialogue is a received literary composition \\
  'Qi hired retainers are no match for Wei's wuzu; Wei's wuzu are no match for Qin's ruishi' -> counter chain of unit types & Xunzi, 'Yi bing' & Xunzi 15 & 据 & The comparative sentence is quoted; the numeric combat coefficients are our conversion \\
  Focus-tree and law nodes carry their sources inside the game (e.g. 'Agronomy and War' quotes the Shang jun shu) & node 'cite' fields in scripts/focus\_tree.gd and scripts/statecraft.gd & scripts/focus\_tree.gd & 推断 & The quoted sources are real; the node effects are game balance \\
  Province anchors use historically attested cities with real coordinates & Tan Qixiang, Zhongguo lishi ditu ji, Warring States sheets & data/provinces.json & 推断 & Some city locations are scholarly conjectures; coordinates digitized by eye from the atlas \\
  Map geometry: province polygons & coastline and mountain polylines & self-drawn vector data rendered by MapView & scripts/map\_view.gd & 设 \\
  Six-stat values for 190 generals & battle records weighted with Shiji verdicts & data/generals.json & 推断 & Derived numbers; rationale recorded but the scale is ours \\
  Formation system (line/wedge/square/arrow) from ancient military practice & Zhoubing texts; game balance & scripts/warfare.gd & 设 & Named formations are attested; the numeric modifiers are invented \\
  King You agent with memory whitelist and one-action JSON (web prototype only) & Park et al. 2023 memory-stream design; Shiji 'Zhou benji' events & prototype sim\_layer.js & 推断 & Built and verified in the prototype; cut from the shipped Godot build at the scope freeze \\
  Round economy & upkeep and combat resolution formula & game rules engine & scripts/warfare.gd + scripts/main.gd & 设 \\
  771 BCE is not the undisputed year of the fall of Western Zhou & Zhang 2023 re-dating of the Zhou accession chronology & Zhang 2023 pp. 18-19 & 据 & Presented as an optional intelligence layer \\
  \bottomrule
\end{tabular}

  \caption{Claims and design elements with locators, evidence tiers
  (\emph{据} documented, \emph{推断} inferred, \emph{设} invented), and limits.}
  \label{tab:evidence}
\end{table}

\section{Verification and Playtesting}

Verification is layered so failure is cheap and diagnosis precise. The first
layer is seventeen headless Godot self-check scripts run against the real
engine binary---map connectivity, campaign flow, city capture, corps
merging, diplomacy and mediation, negotiation, officials, the schools,
focus trees, tactical battles, passage and fortress rules, save--load,
endings, defeat, and a multi-turn soak test---all required to PASS before any
commit. The second layer is the exported web build checked end to end: a
Playwright smoke suite loads both deployed URLs (Cloudflare Workers and
GitHub Pages), asserts every critical asset (\texttt{boot.html},
\texttt{index3.js}, \texttt{index3.wasm}, \texttt{index3.pck}) returns~200,
waits through the software-GL boot, and requires zero console errors and a
rendered canvas. The web prototype additionally carries a WCAG colour-contrast
guard and a turn-loop interaction test, added in response to playtester
feedback.

Two external testers used the build on their own devices. Their feedback was
practical rather than dramatic, and both issues were fixed and locked in. The
first complained that font colours and background were harsh on the eyes
after a few minutes of play; the palette was reworked to soften contrast
against the dark panels, and contrast checks now guard the theme so the harsh
combination cannot be reintroduced. The second found the controls
cumbersome---too many clicks to move armies and confirm each turn; the
movement and turn-confirmation flow was streamlined into fewer, combined
actions, with keyboard shortcuts (\emph{Space} to end the turn) promoted to
first-class citizens and covered by regression tests. Both fixes exist as
permanent tests, so the complaints cannot silently return.

\section{Post-Mortem}
\label{sec:postmortem}

\begin{enumerate}
  \item \textbf{What worked.} Converting the scholarly literature into design
  constraints early---Harnes for the pedagogy, Park for the agent
  architecture, Zhang for the chronology---meant every later argument could
  be settled by pointing at a page number instead of an intuition. The
  single-source-of-truth data pipeline paid for itself: the Godot build
  consumes exported JSON, so numbers cannot fork between the prototype and
  the shipped game. Headless engine tests caught map and rule regressions
  before any human saw them, and the streaming, edge-compressed export made
  a 35-megabyte WebGL game tolerable to load.
  \item \textbf{What failed.} The prototype King You agent leaked future
  history under user provocation; input-side prompting alone proved worthless
  against deliberate attack, and the fix was a world-state validator. The
  first shipping interface was itself a failure found by testers: harsh
  colours and cumbersome controls, which the initial build had treated as
  cosmetic rather than as design defects. The generated portraits, though
  archaeologically grounded in costume, drift in artistic style across the
  set---an unsolved consistency problem.
  \item \textbf{What changed.} The original plan was an open sandbox with a
  scripted tutorial campaign and a conversational King You; the scope freeze
  shipped instead the deeper, more verifiable core---focus trees, statecraft,
  diplomacy, formations---on the Godot build, and left the tutorial and agent
  in the prototype. First-person play of King You had already been dropped in
  favour of treating him as a constrained \emph{other}, following the
  presentism warning in the literature. Naval movement and deep coalition
  diplomacy were deferred, explicitly.
  \item \textbf{Limits and AI tools.} The claim that the knowledge-as-asset
  loop \emph{improves learning} remains a design hypothesis: no pre/post
  measurement has been run, and nothing in the literature licenses the claim
  that games beat films in general. The map's Voronoi geometry does not hug
  the passes and rivers that mattered, and the 87-province economy has not
  been stress-tested past thirty-five turns. The build used Claude and
  GPT-class models for data extraction and code assistance under human review
  of every citation, and image models for the portrait set; a local LLM
  served the prototype agent.
\end{enumerate}

\singlespacing
\printbibliography[title={Bibliography}]

\end{document}
